% Tipo de documento. En este caso es un art�culo, para folios A4, tama�o de la fuente 11pt y con p�gina separada para el t�tulo
\documentclass[a4paper,11pt,titlepage]{article}

% Carga de paquetes necesarios. OrdenesArticle es un paquete personalizado
\usepackage[spanish]{babel} 
\RequirePackage[T1]{fontenc}
\RequirePackage[ansinew]{inputenx} 
\usepackage[spanish,cap,cont,title,fancy]{OrdenesArticle}
\usepackage{lmodern}
\usepackage{array}
\usepackage{graphicx}
\usepackage{hyperref}
\usepackage{pifont}
\usepackage{listings}
\usepackage[usenames,dvipsnames]{color}
\usepackage{colortbl}
\usepackage{color}
\usepackage{ifthen}
\usepackage{longtable}
\hypersetup{bookmarksopen,bookmarksopenlevel=4,linktocpage,colorlinks,urlcolor=blue,citecolor=blue,
						linkcolor=blue,filecolor=blue,pdfnewwindow,
						pdftitle={Aplicaciones para m�viles con SS.OO Symbian},
						pdfauthor={Juan Andrada Romero, Juan Gallardo Casero, Jos� Domingo L�pez L�pez},
						pdfsubject={Ingenier�a del Software II}}


% Macro para definir una lista personalizada 
\newenvironment{milista}%
{\begin{list}{\textbullet}%
{\settowidth{\labelwidth}{\textbullet} \setlength{\leftmargin}{\dimexpr\labelsep+\labelwidth+5pt}
\setlength{\itemsep}{\dimexpr 0.5ex plus 0.25ex minus 0.25ex}
\setlength{\parsep}{\itemsep}
\setlength{\partopsep}{\itemsep}
\addtolength{\topsep}{-7.5pt}
}}%
{\end{list}}

% Macro para insertar una imagen
%       Uso: \imagen{nombreFichero}{Factor escala}{Caption (leyenda)}{Label (identificador para referenciarla)}
% -------------------------------------------------------------------------------------------------------------
\def\imagen#1#2#3#4{
 \begin{figure}[h]
 \begin{center}
   \scalebox{#2}{\includegraphics{#1}}
 \caption {#3}
 \label{#4}
 \end{center}
 \end{figure}
}

% Definicion del "listings" para el lenguaje Java
\lstdefinelanguage{Java}
{
 morecomment = [l]{//}, 
 morecomment = [l]{///},
 morecomment = [s]{/*}{*/},
 morestring = [b]", 
 sensitive = true,
 morekeywords = {package, static, while, switch, break, line, void, String, Object, int, Integer, instanceof, else, if, for, private, return, new, public, class, import, int, boolean, true, false, extends, final, super, protected, abstract, this, do, float, double, null, try, catch, implements}
}

% Configuraci�n del c�digo incluido para el lenguaje Java
\lstset{
  language=Java,
  basicstyle=\footnotesize,
  backgroundcolor=\color{white},
  showspaces=false,
  showstringspaces=false,
  showtabs=false,
  frame=single,
  tabsize=2,
  captionpos=b,
  breaklines=true,
  breakatwhitespace=false,
  escapeinside={\%},
  keywordstyle = \color [rgb]{0,0,1},
  commentstyle = \color [rgb]{0.133,0.545,0.133},
  stringstyle = \color [rgb]{0.627,0.126,0.941}
}

\begin{document}

% En las p�ginas de portada e �ndices, no hay encabezado ni pie de p�gina
\pagestyle{empty} 

% Se incluye la portada
\thispagestyle{empty}
\begin{center}
  {\LARGE UNIVERSIDAD DE CASTILLA-LA MANCHA} \\
  \bigskip
  {\Large ESCUELA SUPERIOR DE INFORM�TICA} \\
  \vspace{28mm}
  \includegraphics[scale=0.45, keepaspectratio]{./imagenes/esi_bw.png} \\
  \vspace{30mm}
  {\LARGE \textbf{ Ingenier�a del Software II}} \\
  \vspace{10mm}
  {\large \textsf{\textsc{SSCA - Sistema de Salud de una Comunidad}}} \\
  {\large \textsf{\textsc{Aut�noma}}} \\
  \vspace{20mm}
  {\large Juan Andrada Romero} \\
  {\large Juan Gallardo Casero} \\
  {\large Jose Domingo L�pez L�pez} \\
  \vspace{9mm}
  {\large \today}
\end{center}
\clearpage

% Texto del reverso de la portada
%%\mbox{}
%%\vspace{18cm}
%%\begin{small}
% Se ajusta la separaci�n entre p�rrafos
%%\parskip=10pt 

%%\copyright~ 2008/2009 Juan Andrada Romero. Universidad de Castilla La Mancha, Escuela Superior de Inform�tica de Ciudad Real.

%%Se permite la modificaci�n, copia y distribuci�n de este documento, seg�n la licencia de documentaci�n GNU (\url{http://www.gnu.org}).

%%Este documento fue compuesto con \LaTeX{}. Im�genes generadas con Power Point y Gimp.
%%\end{small}

%%\newpage

% En las p�ginas de �ndices y prefacio, se utiliza numeraci�n romana
\pagenumbering{Roman}

% Se crea el �ndice
\tableofcontents
% Se pasa p�gina y se a�ade el �ndice de figuras al �ndice principal
%\clearpage\phantomsection
%\addcontentsline{toc}{section}{\listfigurename}
% Se crea el �ndice de figuras
%\listoffigures
% Si se quiere crear un �ndice de tablas se pondr�a: \listoftables

\newpage

% Se ajusta la separaci�n entre p�rrafos
\parskip=10pt

% Se a�ade el prefacio al �ndice
%\clearpage\phantomsection
%\addcontentsline{toc}{section}{Prefacio}

% Comienza el contenido del documento. Se utilizan n�meros ar�bigos y el encabezado y pie de p�gina personalizado
\pagenumbering{arabic}
\pagestyle{fancy}

%Introduccion
\section{Introducci�n}

\subsection{Qu� es Symbian}

Symbian es un sistema operativo (conocido como \textit{Symbian OS}) utilizado en dispositivos m�viles, como PDAs o tel�fonos m�viles. Debido a las limitaciones de estos dispositivos m�viles, Symbian est� dise�ado para residir en un espacio muy peque�o, para utilizar din�micamente los recursos de memoria, para administrar eficientemente la energ�a y para soportar en tiempo real los protocolos de comunicaci�n y telefon�a.
T�cnicamente, el sistema operativo Symbian es una colecci�n compacta de c�digo ejecutable y de otros archivos, siendo la mayor�a de ellos bibliotecas que se cargan din�micamente (DLLs) y archivos de configuraci�n, de im�genes y de tipograf�a. Symbian se almacena, generalmente, en un circuito flash dentro del dispositivo m�vil, para que, gracias a esto, se pueda conservar informaci�n aunque el sistema no posea carga el�ctrica en la bater�a y para permitir la re-programaci�n, sin necesidad de separar este circuito de los dem�s circuitos.

Las aplicaciones compatibles con Symbian se desarrollan a partir de lenguajes de programaci�n orientados a objetos, como C++, Java (con sus variantes como PJava, J2ME, etc.), Visual Basic etc. Cada aplicaci�n desarrollada se puede instalar en la memoria flash del dispositivo con Symbian utilizando un PC, una interfaz USB o firewire, dependiendo del modelo de dispositivo y del cable correspondiente. 

Para terminar, Symbian permite actualizarse, tarea que puede llevar a cabo el propio usuario del dispositivo m�vil a trav�s de los sitios en Internet de los fabricantes de los dispositivos m�viles o bien  al obtener una tarjeta de memoria flash de los distribuidores autorizados, la cu�l contiene la nueva versi�n de Symbian. 

\subsubsection{Breve historia de Symbian}

Symbian OS es una evoluci�n de la plataforma de 32 bits \textit{Psion EPOC} que se utiliz� por primera vez en 1997 en el Psion Series 5 PDA. La primera versi�n de EPOC en el Psion Series 5 se conoce como ''EPOC Release 1'', aunque este n�mero de versi�n iba incrementando pr�cticamente cada a�o, debido a las continuas revisiones que se hac�an de esta primera versi�n.

En 1998, Symbian fue creada como una \textit{spin-off} del \textit{Psion Software Ltd}, co-propiedad de Psion, Nokia, Ericsson y Motorola. El prop�sito de este \textit{spin-off} fue desarrollar una plataforma de software avanzado para unos nuevos modelos de tel�fonos m�viles, llamados tel�fonos inteligentes, que combinan telefon�a y capacidad de c�mputo. Esto represent� un cambio radical en la plataforma de EPOC, que, un a�o m�s tarde, se convirti� oficialmente en la plataforma \textit{Symbian OS}, que es actualmente co-propiedad de diferentes empresas, tales como Motorola, Sony Ericsson y Nokia, entre otros, conocidas como \textbf{Symbian Foundation}. Adem�s, la plataforma Symbian ha sido liberada y en la actualidad es OpenSource.

En cuanto a las versiones de este sistema operativo, el primer lanzamiento de la plataforma Symbian fue la versi�n Symbian OS 6.0, utilizada en el tel�fono \textit{Nokia 9120 Communicator}, en el a�o 2001, evolucionando hasta la versi�n actual, la 9.5. De este modo, aunque se hallan a�adido nuevas funcionalidades al sistema operativo y se haya modificado su API, la arquitectura del sistema operativo Symbian y los paradigmas de programaci�n utilizados se mantienen cercanos a la plataforma original de EPOC de Psion, por lo que no supone mucho esfuerzo migrar de una versi�n a otra.


\subsection{Caracter�sticas t�cnicas de Symbian}


\subsection{Marcas y dispositivos que usan Symbian}

% Mira esta pagina que vienen modelos de m�viles con Symbian: http://aplicaciones-symbian.com/que-es-symbian/

% Esta p�gina est� muy bien para consultar todas las versiones de Symbian y modelos de cada version: http://pdadb.net/index.php?m=os

\clearpage
% Desarrollo
\section{Desarrollo en Symbian}

\subsection{Versiones de Symbian}

En los siguientes apartados se van a comentar las diferentes versiones que existen del sistema operativo Symbian, destacando sus caracter�sticas m�s importantes.

\paragraph{Symbian OS 6.0 y 6.1}

La versi�n 6.0 es la primera versi�n comercial de este sistema operativo, utilizada en el tel�fono m�vil \textit{Nokia 9120 Communicator} en el a�o 2001. A partir de esta versi�n, el sistema operativo de Psion pasa a llamarse Symbian y se encamina al desarrollo para tel�fonos inteligentes \textit{smartphones}. 

En esta versi�n se incluye el soporte para conectividad Bluetooth y el soporte para la ejecuci�n del conjunto de instrucciones de tipo ARM (\textit{Advanced RISC Machine}). Adem�s, se intenta el desarrollo de interfaces de usuario separadas para los \textit{smartphones} y para los dispositivos m�s convencionales, obteniendo dos referentes: la interfaz \textit{Quartz} y la \textit{Crystal}. Sin embargo, aunque se hab�a intentado un desarrollo com�n, se produjo una clara divisi�n, pasando Nokia a utilizar la primera de ellas y Sony Ericsson la segunda. \\
\indent Con esta separaci�n, se comienza a hablar de la serie de dispositivos \textbf{Series 80 1st Edition}, que son los pertenecientes a Nokia, y la serie \textit{UIQ 1.0 y UIQ 1.1}, pertenecientes a Sony Ericsson y Motorola, diferenci�ndose en el tipo de interfaz de usuario, aunque la plataforma y versi�n de Symbian utilizada es la misma en ambos.

En cuanto a la versi�n 6.1, se incorpora el soporte para una c�mara VGA con resoluci�n de 0.3 Mpx, utiliz�ndose por primera vez en el Nokia 7650, perteneciente a la serie \textbf{Serie 601st Edition}.

\paragraph{Symbian OS 7.0}

Esta versi�n se publica en el a�o 2003, a�adiendo soporte para IPv6, para la tecnolog�a EDGE (\textit{Enhanced Data Rates for GSM}) y para Java, siguiendo el est�ndar Java ME. Esta versi�n se utiliza con cualquier tipo de interfaz de usuario, incluyendo las series \textbf{Serie 80 2nd Edition}, \textbf{Serie 90 1st y 2nd Edition}, \textbf{Serie 60 2n Edition}, \textbf{UIQ 2.0} y \textbf{UIQ 2.1}.

Por otra parte, �sta es la primera versi�n de Symbian con un problema de seguridad grave, ya que es infectada por el gusano llamado \textit{Cabir}, capaz de transmitirse a otros m�viles con Symbian OS a trav�s del Bluetooth.

\paragraph{Symbian OS 8.0}

La versi�n 8.0, lanzada en el a�o 2004, cuenta con dos kernel del sistema operativos diferentes, el EKA1 y el EKA2, aunque �ste ultimo no se incorpora hasta la versi�n 8.1b. En esta versi�n se mejoran las capacidades multimedia del sistema operativo, teniendo mas rendimiento en aplicaciones que utilizan \textit{streaming} e incorporando OpenGL con gr�ficos vectoriales; se permite la administraci�n remota del dispositivo; se soporta CDMA (Code Division Multiple Access) y la tecnolog�a 3G y, por �ltimo, se aumentan los perif�ricos de entrada y salida, dando soporte a sintonizadores digitales de TV (DVB-H), reconocimiento de huellas, etc.

Esta versi�n se incluye en las serie \textbf{Series 60 2nd Edition 2.6}.

\paragraph{Symbian OS 8.1}

\paragraph{Symbian OS 9.0}

\paragraph{Symbian OS 9.0}

\paragraph{Symbian OS 9.1}

\paragraph{Symbian OS 9.2}

\paragraph{Symbian OS 9.3}

\paragraph{Symbian OS 9.4}

\paragraph{Symbian OS 9.5}

\subsection{Entornos de desarrollo, API y SDK de Symbian}


\subsection{Firma y publicaci�n de aplicaciones}
\clearpage

% Secci�n de ap�ndices
% Se a�ade el titulo de "Ap�ndices" al �ndice
%\addcontentsline{toc}{section}{Ap�ndices}

%\input{Apendices.tex}

% A�adimos la bibliografia al �ndice
\clearpage\phantomsection
\addcontentsline{toc}{section}{Referencias}

% Bibliograf�a. En este caso se usa BibTeX
\bibliographystyle{plain}
\pagestyle{plain} 
\bibliography{Bibliografia}

\end{document}
